\subsection{A Gaussian-relative bound for the complete source}

Set \(p(u)=Tu^2+9u-\pi e^{4u}\), \(W=\sqrt{\kappa}(U-L)\) and
\(\alpha=2S/(3\kappa^{3/2})\). In the \(w\)-coordinate the unnormalized
density, with a convenient constant, is
\[
 F(w)=\frac{e^{p(u)-p(L)}}{\sqrt{2\pi}}\,b(u),\qquad
 u=L+\frac{w}{\sqrt\kappa},\qquad
 b(u)=\sum_{n\ge1}\left(n^4-\frac{3n^2e^{-4u}}{2\pi}\right)
                         e^{-\pi(n^2-1)e^{4u}} .
\]
Extend \(F\) by zero for \(w<-L\sqrt\kappa\) and put
\(\varphi(w)=e^{-w^2/2}/\sqrt{2\pi}\).

\begin{lemma}[Centered source on fixed and growing disks]\label{lem:saddle}
For each fixed \(r>0\), and also for \(r=\tfrac12\sqrt{\log\kappa}\), the
logarithm from zero satisfies, for all sufficiently large \(T\),
\begin{equation}
 \sup_{|\zeta|\le r}
 \left|K_0(\ell\zeta)+\frac{\zeta^2}{2}-\ii\alpha\zeta^3\right|
 \le E,\qquad
 E=O\!\left(\frac{(1+r^6)e^{r^2/2}}{\kappa}\right).
 \label{eq:window-log}
\end{equation}
In particular the fixed-disk estimate is
\begin{equation}
 K_0(\ell\zeta)=-\zeta^2/2+\ii\alpha\zeta^3+O_r(\kappa^{-1}).
 \label{eq:complex-saddle}
\end{equation}
The exact mean and third cumulant satisfy
\(a=L-2S/\kappa^2+O(\kappa^{-3/2})\) and
\(\gamma_3=-4S/\kappa^3+O(\kappa^{-5/2})\).
\end{lemma}
\begin{proof}
The following majorants give a finite test as well as the asymptotic
estimate \cite{EZWindow}. Put \(w_0=\kappa^{1/12}\) and require
\(w_0\le\sqrt\kappa/4\), \(L\ge1/4\) and
\[
 c_*=\frac{S/e-2T}{2\kappa}>0,\qquad
 \Delta=\min\{p(L)-p(0),c_*\kappa/16\}>0 .
\]
For \(u\ge L-1/4\), \(p''(u)\le-2\kappa c_*\), giving
\(p(u)-p(L)\le-c_* w^2\). Since \(p''\) decreases, any convex part to the
left of \(L-1/4\) attains its maximum at an endpoint, and the concave
part increases up to \(L\). Thus \(p(u)-p(L)\le-\Delta\) for
\(0\le u<L-1/4\).

For all \(u\ge0\), \(0<b(u)\le b_*=(1-16e^{-3\pi})^{-1}\).
On the central interval \(|w|\le w_0\), use
\[
 v_b=16e^{-3S/(16e)}<1,\qquad b_{\rm err}=\frac{24e}{S}+\frac{v_b}{1-v_b}.
\]
Here \(\pi e^{4u}\ge S/(16e)\) and the first summand differs from one
by at most \(24e/S\). The tail follows from \(n^4\le16^{n-1}\) and
\(n^2-1\ge3(n-1)\). These give \(|b-1|\le b_{\rm err}\).
Since \(p'''(L)=-4S\) and \(|p^{(4)}|\le16eS\) on the central
interval, the exponent is \(-w^2/2-\alpha w^3\) with error at most
\(\beta w^4/\kappa\), where \(\beta=2eS/(3\kappa)\). It follows that
\[
 |F(w)-\varphi(w)(1-\alpha w^3)|
 \le \frac{A}{\kappa}\varphi(w)(1+w^4+w^6),\qquad
 A=e^{\alpha w_0^3+\beta w_0^4/\kappa}
        \max\{\kappa b_{\rm err},\beta,\kappa\alpha^2/2\}.
\]
Indeed changing the exponent costs at most the exponential factor
in \(A\) times the exponent error. The quadratic remainder of
\(e^{-\alpha w^3}\) and the factor \(b-1\) give the other two terms.

Let \(\mathcal F(\zeta)=\int_{\mathbb R}F(w)e^{\ii\zeta w}\dd w\).
The Gaussian model has transform
\(e^{-\zeta^2/2}\{1+\ii\alpha(\zeta^3-3\zeta)\}\). For
\(\zeta=\xi+\ii\eta\), completing the Gaussian square in the central
error under \(e^{-\eta w}\) costs \(e^{\eta^2/2}\). Division by the
modulus of the Gaussian transform leaves \(e^{\xi^2/2}\). The shifted
moment of \(1+w^4+w^6\) is \(19+51\eta^2+16\eta^4+\eta^6\).
Together with both source tails and the Gaussian model tail this gives
the following bound on the relative unnormalized error for \(|\zeta|\le s\):
\begin{align}
 A_s=e^{s^2/2}\Bigg[&
  \frac{A}{\kappa}(19+51s^2+16s^4+s^6)\nonumber\\
 &+\int_{|w|>w_0}e^{s|w|}
    \left(\frac{b_*}{\sqrt{2\pi}}e^{-c_* w^2}
                     +\varphi(w)(1+\alpha|w|^3)\right)\dd w \nonumber\\
 &+\frac{b_*}{\sqrt{2\pi}}L\sqrt\kappa\,e^{-\Delta+sL\sqrt\kappa}\Bigg].
 \label{eq:window-source-test}
\end{align}
Write \(P=\alpha(r^3+3r)\). If \(A_0<1\), normalization by
\(\mathcal F(0)\) adds at most the amount charged in
\[
 \epsilon=\frac{A_r+(1+P)A_0}{1-A_0},\qquad v=P+\epsilon .
\]
For \(v<1\), the normalized relative transform lies in the logarithm
disk about one. Its logarithm differs from \(\ii\alpha(\zeta^3-3\zeta)\)
by at most \(\epsilon+v^2/[2(1-v)]\), by the logarithm series.
Centering subtracts the actual linear coefficient. Cauchy's estimate
bounds the subtraction of the linear term of the residual on
\(|\zeta|\le r\) by one more copy of that residual. This proves the finite
version of \eqref{eq:window-log} with
\begin{equation}
 E=2\left(\epsilon+\frac{v^2}{2(1-v)}\right).
 \label{eq:window-error-test}
\end{equation}

Eventually \(S/\kappa\to1\), \(c_*\to1/(2e)\), \(A=O(1)\), and
\((p(L)-p(0))/\kappa\to\infty\). For either radius in the statement,
\(r=o(w_0)\) absorbs the weighted Gaussian tails into quadratic
decay, and \(rL\sqrt\kappa=o(\kappa)\) leaves exponential decay in
the far-left tail. Thus \(A_0=O(\kappa^{-1})\) and
\(A_r=O((1+r^6)e^{r^2/2}/\kappa)\), including the factor from
division by the Gaussian. These bounds imply \(v=o(1)\) and the
stated error bound. For a fixed radius the first uncentered
logarithmic coefficient gives \(\E_T W=-3\alpha+O(\kappa^{-1})\).
The cubic coefficient of \eqref{eq:complex-saddle} gives
\(\gamma_3=-6\alpha/\kappa^{3/2}+O(\kappa^{-5/2})\).
\end{proof}

\subsection{The count and the position and gap errors}

For the disk in \eqref{eq:window-log} put
\(R=\ell r\), \(M=\alpha r^3+E\), \(b_R=M/6\) and \(Q=R/2-b_R\).
The odd part of \eqref{eq:window-log} gives
\[
 \sup_{|y|\le R}|\Psi(y)+\alpha(y/\ell)^3|\le E,\qquad
 \sup_{|y|\le R}|\Psi(y)|\le M .
\]
If \(M\le R/4\), Theorem~\ref{thm:inverse} supplies both the true and
cubic inverse branches in \(|q|<Q\), since even the full coefficient
majorant on the boundary is \(b_R\). On \(|y|\le R/2\),
Cauchy's estimate gives \(|\Psi'(y)|\le2M/R\le1/2\).
For \(q=q_k\) and \(q'=q_{k+1}\), define the common residual
\(\mathcal R(q)=Y(q)-Y_3(q)\). In physical units the errors are
\begin{equation}
 x_k-Y_3(q)/a=\mathcal R(q)/a,\qquad
 g_k-\{Y_3(q')-Y_3(q)\}/a
      =\{\mathcal R(q')-\mathcal R(q)\}/a .
 \label{eq:shared-residual}
\end{equation}

\Needspace{8\baselineskip}
\begin{theorem}[Growing counted window]\label{thm:growing-window}
With \(r=\tfrac12\sqrt{\log\kappa}\), the tests above, including
\(M\le R/4\), hold for all sufficiently large \(T\). There are exactly
\begin{equation}
 K(T)=\left\lfloor\frac{\ell\sqrt{\log\kappa}}{8\pi}+\frac12\right\rfloor
 \sim\frac{\sqrt{2}}{32\pi}\sqrt T(\log T)^2
 \label{eq:growing-count}
\end{equation}
positive zeros through \(Y(R/4)/a\). All are simple, including a
crossing at the endpoint when it occurs. Uniformly for \(q\le R/4\)
in position and \(q'\le R/4\) in a gap,
\begin{align}
 |x_k-Y_3(q)/a|
 &\le\frac{E}{a}\left(\frac{q}{Q}\right)^5,\nonumber\\
 \left|g_k-\frac{Y_3(q')-Y_3(q)}{a}\right|
 &\le \frac{7\pi E}{aQ}\left(\frac{q'}{Q}\right)^4 .
 \label{eq:growing-errors}
\end{align}
The last zero has height \(x_{K(T)}\sim(\sqrt{2}/8)\sqrt T\log T\)
and \(g_k=(\pi/a)(1+O(\kappa^{-1}))\) throughout the counted window.
\end{theorem}
\begin{proof}
The inverse count gives exactly the levels \(q_k\le R/4\).
The physical height uses \(x_k=q_k/a+O(M/a)\),
\(a\sim L\sim(\log T)/4\) and \(\kappa\sim2T\log T\).
The spacing bound follows from \(Y'=1/(1-\Psi'(Y))\) and
\(M/R=O(\kappa^{-1})\) at the growing radius.

Cauchy's bound for the cubic coefficient gives
\(\lvert c_1\ell^3+\alpha\rvert\le E/r^3\), so \(c_1<0\) eventually.
The analytic function \(I(y)=\Psi(y)-c_1y^3\) vanishes to order five
and has supremum at most \(2E\) on \(|y|\le R\). Divide by \(y^5\)
to obtain \(|I(y)|\le2E(|y|/R)^5\). Subtract the inverse equations:
\[
 \mathcal R(q)\{1-c_1(Y^2+YY_3+Y_3^2)\}=I(Y).
\]
The brace has modulus at least \(1-3M/(4R)\ge13/16\).
Since \(|Y|<R/2\), this gives \(|\mathcal R|\le E\) on \(|q|<Q\);
the difference vanishes to order five. For any holomorphic \(f\) on
\(|q|<Q\), bounded by \(H_*\) and vanishing to order \(m\), the maximum
principle and Cauchy's estimate for \(f/q^m\) give, for \(0<z<Q\),
\[
 |f(z)|\le H_*(z/Q)^m,\qquad
 |f'(z)|\le \frac{H_*}{Q}(z/Q)^{m-1}\left(m+\frac{z}{Q-z}\right).
\]
Here \(Q\ge3R/8\), so the factor in parentheses is at most seven for
\(m=5\) and \(z\le R/4\). Integrate the derivative over \([q,q']\)
and use \eqref{eq:shared-residual}.
\end{proof}

At the growing radius, \(E=O(\kappa^{-7/8}(\log\kappa)^3)\).
At the outer scale the errors in \eqref{eq:growing-errors} are
\(O(\kappa^{-7/8}(\log\kappa)^2)\) for a physical position and
\(O(\kappa^{-11/8}(\log\kappa)^{3/2})\) for a gap in units of \(\pi/a\).

\subsection{Sharper errors on fixed scaled windows}

For fixed \(C>0\), repeat the argument with \(r=4C\).
Then \(Q\ge2C\ell-O_C(\kappa^{-1/2})\), \(E=O_C(\kappa^{-1})\),
and the counted interval \(q_k\le C\ell\) has length
\(\lfloor C\ell/\pi+1/2\rfloor\) and last height asymptotic to
\(C\sqrt\kappa\). Since the odd saddle residual vanishes to order
three, \(|Y-q|=O_C(\kappa^{-1/2}(q/\ell)^3)\), using \(|Y|\le2q\).
Replacing \(Y\) by \(q\) in the saddle cubic then gives
\begin{equation}
 ax_k=q-\alpha(q/\ell)^3+O_C(\kappa^{-1}(q/\ell)^3),\qquad
 |ax_k-Y_3(q)|=O_C(\kappa^{-1}(q/\ell)^5).
 \label{eq:distort}
\end{equation}
For \(q'\le C\ell\), the corresponding normalized gap error is
\begin{equation}
 \frac{|\mathcal R(q')-\mathcal R(q)|}{\pi}
   =O_C\!\left(\frac{1}{\kappa\ell}
                    \left(\frac{q'}{\ell}\right)^4\right).
 \label{eq:fixed-gap}
\end{equation}

For each fixed \(k\ge2\),
\begin{equation}
 \lim_{T\to\infty}L^3\kappa^2
 \left(1-\frac{x_k(T)}{(2k-1)x_1(T)}\right)
  =\frac{2\pi^2}{3}k(k-1).
 \label{eq:shape}
\end{equation}
In the ratio the common factor \(1/a\) cancels. The leading distortion
comes from \(c_1 q^3\) in \eqref{eq:inverse}, with
\(q_k^2-q_1^2=\pi^2k(k-1)\) and \(a/L,S/\kappa\to1\).
The remaining terms are smaller by \eqref{eq:distort} and
\(Y_3-q-c_1q^3=O_k(c_1^2)\).
For fixed \(k\), the leading shift of the common frequency gives
\[
 x_k(T)=\frac{q_k}{L-2S/\kappa^2}
       +O_k\!\left(\frac{1}{L^2\kappa^{3/2}}\right).
\]
The two-site Jacobi reference has frequency \(\sqrt{\beta_1}\) with
\(\beta_1=\E_T U^2\) and \(\beta_1\beta_2=\Var_T(U^2)\).
Put \(\epsilon_*=\sqrt{\Var_T(U)}/a\) and \(d_J=\beta_2/\beta_1\).
The centered moment estimates at the saddle give
\(d_J=4\epsilon_*^2+O(\epsilon_*^3)\). Since
\(\sqrt{\beta_1}/a=\sqrt{1+\epsilon_*^2}\) and
\(ax_k=q_k+O_k(\epsilon_*^3)\), the squared half-gap satisfies
\[
 \frac{D_k}{\pi^2/(4\beta_1)}=1+\frac{d_J}{4}+O_k(d_J^{3/2})
 \qquad\text{for fixed }k .
\]
This compares the gap with the one from that Jacobi reference
\cite[\S2]{Founding07}.

\subsection{Finite source data and direct phase localization}
\label{sec:inverse-data}

At a specified \(T\), evaluate \eqref{eq:window-source-test} and
\eqref{eq:window-error-test} with enclosed saddle and source quantities,
check the stated inequalities, and retain levels with
\(q_k\le R/4\) using a lower bound for \(R\). The tails in
\eqref{eq:window-source-test} reduce to elementary Gaussian integrals.
The asymptotic theorem leaves the numerical onset in \(T\) unevaluated.

With the exact \(a\), truncate \eqref{eq:inverse} at degree \(2N+1\).
Set \(m=2N+3\) and \(s=q/Q\). Cauchy's coefficients for \(Y-q\)
bound the physical position tail by \((b_R/a)s^m/(1-s^2)\).
For the adjacent difference, differentiate the geometric majorant. At
\(s=q'/Q\) the bound is
\begin{equation}
 \frac{\pi b_R}{aQ}
  \left(\frac{m s^{m-1}}{1-s^2}
                       +\frac{2s^{m+1}}{(1-s^2)^2}\right).
 \label{eq:inverse-tail}
\end{equation}
If the coefficient of \(q^{2j+1}\) has absolute error at most \(d_j\),
add \(a^{-1}\sum_{j=1}^N d_j q^{2j+1}\) in position and
\((\pi/a)\sum_{j=1}^N(2j+1)d_j(q')^{2j}\) in spacing. The polynomial
expressions in \eqref{eq:inverse} propagate the input cumulant intervals
to these coefficient errors.

To compare with direct localization from identical data, take the phase
polynomial \(\Psi_N=\sum_{j=1}^N c_j y^{2j+1}\), \(N\ge1\).
With exact coefficients, Cauchy's estimate on the saddle residual
in \eqref{eq:window-log} gives
\[
 \sup_{|y|\le R/2}|\Psi(y)-\Psi_N(y)|
 \le \delta_N=\tfrac43 E\,2^{-m}.
\]
The truncated polynomial uses the same boundary majorant \(b_R\);
its inverse \(Y_N\) lies in the same disk through the retained levels.
Subtract the two inverse equations. The true phase perturbation
\(\Psi\) has derivative bounded by \(1/2\) on the convex disk
\(|y|\le R/2\), so \(|Y-Y_N|\le2\delta_N\) on \(|q|<Q\).
The residual vanishes to order \(m\) at zero. The position and gap
bounds from the quotient and its derivative are respectively
\[
 \frac{2\delta_N}{a}(q/Q)^m,\qquad
 \frac{2\pi\delta_N(m+2)}{aQ}(q'/Q)^{m-1}.
\]
On the real interval, values of \(y-\Psi_N(y)-q_k\) below
\(-\delta_N\) and above \(\delta_N\) give signed brackets for the
same true crossings. At \(N=1\), the phase polynomial is inverted by
the closed cubic \(Y_3\).
A direct sign test of the field itself pays the amplitude as well:
on \(|y|\le R/2\) the real exponential factor in the field is at least
\(e^{-r^2/8-E}\) by \eqref{eq:window-log}. Thus the two approaches share
the count and the analytic error for direct use of the same cubic.
The inverse series supplies reusable coefficients for all ranks in the
window.

A common perturbation analysis prices inexact phase coefficients. At
the exact coordinates, coefficient errors \(\varepsilon_j\) cost
\(\sum_j\varepsilon_j(R/2)^{2j+1}\) on the boundary; include any omitted
tail in the total bound \(\delta_*\). If
\(Q_*:=R/2-b_R-\delta_*>R/4\), Rouch\'e's theorem keeps the
perturbed inverse in \(|y|<R/2\) for \(|q|<Q_*\). Subtraction of the
equations again bounds the displacement by \(2\delta_*\).
Cauchy's estimate gives physical adjacent error at most
\(2\pi\delta_*/[a(Q_*-R/4)]\).
For an error \(|\widehat c_1-c_1|\le d_c\) in the cubic alone, take
\(\delta_*=d_c(R/2)^3\). If \(\delta_*\le R/16\) then
\(Q_*\ge5R/16\), and the cubic difference vanishes to order three.
The same quotient bounds give \(O(d_c q^3/a)\) in position and
\(O(\pi d_c(q')^2/a)\) in the gap. Converting with
\(|\widehat a-a|\le d_a\le a/2\) adds
\(O(d_a q/a^2)\) and \(O(\pi d_a/a^2)\), respectively. Indeed the
perturbed inverse is \(O(q)\) with bounded slope on the retained
segment by its strict radius and its zero at the origin. The
coefficient budget must include uncertainty in \(a\) when forming
\(\widehat c_1\). At the outer scale \(q\asymp R\), the orders
\(d_c=O(E/R^3)\), \(d_a=O(aE/R)\) preserve the errors in
\eqref{eq:growing-errors}. Near the origin the respective powers of
\(q\) in these bounds determine the smaller admissible errors.

The analytic envelope adds information beyond any finite list of
moments in the class of unrestricted positive half-line densities.
To see this, add to the positive density of \(U\) a nonnegative
bump translated to \(u=H_*\) with mass \(H_*^{-m-1}\). Its first \(m+1\)
moments, starting with degree zero, are \(O(H_*^{-1})\). Choose \(m+1\)
fixed interior bumps with an invertible moment matrix by taking small
supports about distinct positive points. Subtract their linear
combination with matching moment vector. The coefficients tend to
zero, so positivity and all the specified moments are preserved.
The exponential moment at each fixed positive rate tends to infinity.
The inverse remainder proved here therefore uses the analytic envelope
as an additional input. The individual contour formula \eqref{eq:loop}
also requires a whole-circle error and real clearance in addition to
the derivatives at its center.

These estimates use the normalized complete theta expansion for the
source envelope and the literal cosine for the zero count. Euler-factor
changes alter the complete source coefficients. Weighted-square and
finite-prime perturbations change the unit weights and full
reflection. Reflected, smooth, and primitive-compatible controls lack
the canonical square expansion. Atomic threshold controls have a
different support, and finite-profile and finite-tail variants change
the complete source even when selected observations agree. The
normalized endpoint perturbations change Dirichlet series arithmetic.
A changed observation kernel loses the cosine factorization used here
\cite{EZWindow}. A control with its own verified envelope
can use the same inverse argument on the corresponding disk.
The gain for the theta field is a larger counted family and a
derivative bound for the shared gap error at one time. The count ends
at the proved height.
